\section{Introduction}
\label{sec:introduction}

The Wolfram Physics Project~\cite{wolfram2020project} models the universe as an abstract rewrite system, in which physics is given by the causal structure of the rewrites rather than by the data being rewritten. The states are hypergraphs, which have no built-in notion of space: distance is path length through the edges. Rewriting produces three structures: the multiway evolution graph of all reachable states and the events between them; the causal graph, of dependencies between events; and the branchial graph, of pairs of events that consume a common edge of the same state.

Exhaustive multiway evolution is expensive. Each step multiplies the number of histories by the frontier's match count, and a rule that adds edges gives each state more matches than its parent, so the count grows super-exponentially: on the canonical Wolfram rule the ratio between consecutive depths rises from $1.7$ at depth two to $11.7$ at depth seven. Every step matches the rule against each new state, canonicalizes each new state to detect isomorphic duplicates, and records events and their causal and branchial relations. The existing Wolfram Language implementation takes $\AuthorityDeepSeconds$\,s to reach depth \SpeedupHighDepth{} of the corpus's \texttt{wolfram-2to4} rule; this engine takes $\CoreDeepMs$\,ms for the same evolution (Section~\ref{sec:benchmarks}).

The engine provides:

\begin{enumerate}
    \item \textbf{Exact canonicalization:} Every state is canonicalized by McKay-style individualization--refinement (IR)~\cite{mckay2014practical}, and the canonical hash is the deduplication key, so two states are identified if and only if they are isomorphic. Proposition~\ref{prop:wlceiling} bounds the saving from any cheaper filter in front of IR by the ratio of canonical classes to raw states, which falls toward zero with depth.

    \item \textbf{Lock-free concurrency:} The engine holds no lock. Its hash maps, segmented arrays, lists and work-stealing deques are non-blocking, and an idle worker waits on an atomic counter through \texttt{futex}~\cite{franke2002fuss} on Linux, \texttt{WaitOnAddress} on Windows or \texttt{os\_sync\_wait\_on\_address} on macOS, using no cycles until new work arrives.

    \item \textbf{GPU acceleration:} A CUDA kernel that stays resident for the whole evolution runs matching, rewriting and canonicalization, so no step returns to the host. It is faster than the CPU when the rule set gives it enough independent work (Section~\ref{sec:benchmarks}).

    \item \textbf{Configurable equivalence:} State and event identity are chosen independently. State identity is \textsc{None}, \textsc{Automatic} (a content hash) or \textsc{Full} (exact IR); event identity is built from a small set of signature keys.

    \item \textbf{Exploration strategies:} Quotient exploration (Section~\ref{sec:quotient}) expands one state per isomorphism class and still reconstructs every raw event; sampling with per-rule rates and a per-state match cap (Section~\ref{sec:sampling}) reaches depths where exhaustive evolution cannot; and the causal graph's transitive reduction (Section~\ref{sec:tr}) is maintained as edges are added.

    \item \textbf{Determinism:} The state set, event set and causal and branchial relations depend only on the inputs, not on the thread count or the scheduling order (Section~\ref{sec:determinism}).
\end{enumerate}

Section~\ref{sec:background} reviews hypergraph rewriting and related work; Section~\ref{sec:canonicalization} describes canonicalization; Section~\ref{sec:pattern-matching} the pattern matcher; Section~\ref{sec:architecture} the system architecture; Section~\ref{sec:gpu} the GPU implementation; Section~\ref{sec:benchmarks} the evaluation; and Section~\ref{sec:verification} the verification of the concurrent code. Section~\ref{sec:conclusion} concludes.

% ----------------------------------------------------------------------------
% BACKGROUND
% ----------------------------------------------------------------------------
